% Fibonacci-Pell nearest gaps: an all-exponent classification and
% quadratic-unit orbit rigidity
%
% Authors:
%   Dr. Denys Dutykh (Mathematics Department, Khalifa University of Science
%   and Technology, Abu Dhabi, UAE)
%   Prof. Laurent Vuillon (Univ. Savoie Mont Blanc, CNRS, LAMA, Chambery,
%   France)
%
\section{Scope, connections, and the next obstruction}
\label{sec:outlook}

The first arithmetic conclusion of this paper is complete on its stated
domain: for every $q\geqslant1$, both signs, and every positive exponent,
the strict Fibonacci--Pell norm--content system has exactly the two rows in
\cref{eq:two-solutions}, and the proof is unbounded in $q$, closing only
after the analytic bounds have reduced both parity branches to exhaustive
finite domains. The second conclusion is likewise complete:
\cref{thm:complete-fibonacci-orbits} gives all orbit hits for both signed
Fibonacci cores, every positive Pell anchor, and every nonnegative time.

The surrounding results explain why these classifications are rigid. At the
Archimedean place, candidate windows are abundant and their exact densities
are Haar measures of interval intersections on a phase circle; at split
finite places, compatible Fibonacci and companion ranks can synchronize to
arbitrary prescribed depth. Neither phenomenon sees the oriented primitive
unit direction, whereas the global norm--content equation sees both
coefficient coordinates at once, and it is this orientation that enables
the primitive rectangle and the logarithmic-form argument. The separation
is therefore not local versus computational, but measure-theoretic
abundance versus global Diophantine rigidity.

The algebraic structures exposed here have precise roles: the norm torus,
its negative-norm torsor, the integral unit group, and the Lorentzian
symmetric-square representation all act directly in the proofs, while the
biquadratic field $\Q(\sqrt2,\sqrt5)$ supplies the two conjugations used in
the norm and height arguments. Further frameworks would become relevant
only after an explicit map or invariant connects them to the classified
system.

There is, however, a genuine algebraic-geometric object behind the parent
problem. Over $\Z$, the Markoff equation defines the affine cubic scheme

\begin{equation}
 \mathcal M
 =\operatorname{Spec}
 \Z[a,b,c]/(a^2+b^2+c^2-3abc),
 \label{eq:markoff-scheme}
\end{equation}

and the Vieta involutions, such as

\begin{equation}
 (a,b,c)\longmapsto(a,b,3ab-c),
 \label{eq:vieta-involution}
\end{equation}

are integral automorphisms of this cubic. After the scaling
$(x,y,z)=(3a,3b,3c)$, its equation becomes
$x^2+y^2+z^2-xyz=0$, the Fricke trace surface with commutator trace $-2$.
This is also a concrete algebraic-topological bridge: the fundamental group
of a one-holed torus is the free group $F_2$, and the three coordinates are
traces of two generators and their product in an
$\mathrm{SL}_2$-representation. Goldman shows that the reflection

\[
 (x,y,z)\longmapsto(x,y,xy-z)
\]

comes from the extended mapping-class group of the surface
\cite[Sections~2.1--2.2 and Appendix~A-3]{Goldman2003}, and scaling back
gives exactly \cref{eq:vieta-involution}. On the discrete side, taking
positive Markoff triples as vertices and Vieta moves as edges gives the
mutation graph and its familiar descent tree; Christoffel words and integral
necklaces encode its arithmetic vertices
\cite{Reutenauer2026,Fisac2025}. Topology thus explains the symmetry, graph
theory records the orbits, and combinatorics labels the vertices. These
connections are exact but are not hidden inputs: we construct no morphism
from a hypothetical collision locus on $\mathcal M$ to the Fibonacci
nearest-gap system.

The most immediate arithmetic extension starts instead from
\cref{thm:general-orbit-gap}. For every noncanonical positive seed, an orbit
hit is equivalent to a strict norm--content gap for the seed-dependent
quadratic integer $A_z$. The Fibonacci cores are one coefficient locus, and
on it we obtain a complete classification; noncanonical cores with
$\mathcal Q(z)^2>1$ lead to more general continuant-dependent coefficients.
The next problem is to reconstruct enough primitive geometry from those
coefficients to retain effective heights without assuming the very collision
one hopes to exclude. It belongs in a separate paper, since it needs new
arithmetic rather than more exposition around the two-row theorem.

The moving-coefficient logarithmic form should also be understood at the
level of mechanism. Grouping all terms of the same dominant growth before
applying a logarithmic-form estimate is established practice in exponential
Diophantine equations
\cite{DdamuliraLucaRakotomalala2016,BravoDasHerreraSaunders2026}; what is
specific here is the coefficient \cref{eq:moving-coefficient}, its exact
full norm \cref{eq:A-norm}, and the explicit reduction turning those
identities into a classification. We claim no priority for the strategy.

The Markoff-facing bottleneck is therefore sharp. Downstream of the
Fibonacci-core hypothesis, the gaps and all positive-anchor orbits are
classified; upstream, one must prove that a repeated largest Markoff number
yields one of these cores, find a competing core, or classify the gaps with
$\mathcal Q(z)^2>1$.

One tempting fourth route meets a sharp obstruction. The map $\mathsf S$ is
an integral square root of the Pell arm's monodromy Pisot unit
$\beta(2)=\lambda^2$, and \cref{thm:half-step-criterion} shows that an
analogous half-step on the branch through $m_0$ can exist only when
$3m_0-2$ is a perfect square.
One displayed ray fixing $m_0=34$ shows that the arithmetic condition alone
does not guarantee compatibility: its two-square data admit no common
integral half-step, though other fixed-$34$ rays are not classified here.
Transport to another branch would need both the square condition and a
compatible action on that ray's two-square sequence, and the branch
recurrence alone supplies neither that compatibility nor its failure.

The frontier is precise: local abundance versus global rigidity. Progress on
Markoff uniqueness requires crossing the missing upstream arrow in
\cref{fig:logic}.
